\documentclass[12pt, a4paper]{article}

\usepackage[utf8]{inputenc}
\usepackage[T2A]{fontenc}
\usepackage[ukrainian]{babel}
\usepackage{amsmath, amssymb, amsthm}
\usepackage{mathtools}
\usepackage{geometry}
\usepackage{hyperref}
\usepackage{booktabs}
\usepackage{array}
\usepackage{enumitem}
\usepackage{algorithm}
\usepackage{algpseudocode}
\usepackage{lmodern}
\usepackage{tabularx}

\geometry{margin=2.5cm}

\hypersetup{
    colorlinks=true,
    linkcolor=blue,
    urlcolor=blue,
    citecolor=blue
}

\theoremstyle{definition}
\newtheorem{definition}{Визначення}[section]
\newtheorem{remark}{Зауваження}[section]

\title{%
  \textbf{Математична постановка задачі оптимізації}\\[0.5em]
  \large DeliveryIQ --- Оптимізатор маршрутів
}
\author{}
\date{}

\begin{document}

\maketitle
\tableofcontents
\newpage

% ─────────────────────────────────────────────────────────────────────────────
\section{Загальний огляд}
% ─────────────────────────────────────────────────────────────────────────────

DeliveryIQ розв'язує дві пов'язані комбінаторні задачі оптимізації залежно
від кількості транспортних засобів у парку:

\begin{enumerate}
  \item \textbf{Режим одного транспортного засобу} (розмір парку $= 1$):
        \emph{Задача комівояжера} (TSP) --- знайти найкоротший маршрут з
        поверненням до депо, що відвідує всі точки доставки рівно по одному разу.

  \item \textbf{Режим кількох транспортних засобів} (розмір парку $\geq 2$):
        \emph{Задача маршрутизації транспортних засобів з обмеженнями
        вантажопідйомності} (CVRP) --- розподілити точки між різнорідним парком
        транспортних засобів з урахуванням модального типу та максимальної
        кількості зупинок для кожного, а потім розв'язати TSP для кожного
        призначеного кластеру.
\end{enumerate}

Обидві задачі функціонують на \emph{зваженому орієнтованому графі}, отриманому
з OpenStreetMap (OSM) за допомогою бібліотеки OSMnx. Часи переміщення вздовж
ребер графа кодують правила модальної доступності (автомобіль / велосипед /
пішки) і є цільовою функцією оптимізації.

% ─────────────────────────────────────────────────────────────────────────────
\section{Модель графа дорожньої мережі}
% ─────────────────────────────────────────────────────────────────────────────

\subsection{Орієнтований зважений граф}

Вулична мережа моделюється як орієнтований мультиграф

\[
  G = (V,\, E,\, w),
\]

\noindent де:
\begin{itemize}
  \item $V$ --- множина OSM \emph{вузлів} (перехрестя та кінці доріг),
  \item $E \subseteq V \times V$ --- множина орієнтованих \emph{ребер}
        (ділянки доріг),
  \item $w : E \to \mathbb{R}_{>0} \cup \{\mathcal{P}\}$ ---
        \emph{функція вагових коефіцієнтів часу переміщення}, що відображає
        кожне орієнтоване ребро на час проходження у секундах або на
        штрафний сигнал $\mathcal{P} = 10^9$ (непрохідне ребро).
\end{itemize}

\subsection{Найбільша сильно зв'язна компонента (НСЗК)}

Після завантаження граф обрізається до своєї \emph{найбільшої сильно
зв'язної компоненти}:

\[
G^* = \mathop{LSCC}(G),
\]

\noindent що гарантує існування орієнтованого шляху з $u$ до $v$ у межах $G^*$
для кожної пари вузлів $u, v \in V^*$. Усі подальші обчислення використовують
виключно $G^*$.

\subsection{Присвоєння часів переміщення за модальними типами}
\label{sec:modal}

Будуються три незалежні копії $G^*$ --- по одній для кожного транспортного
режиму $m \in \mathcal{M} = \{\texttt{drive},\, \texttt{bike},\, \texttt{walk}\}$.
Ваговий коефіцієнт часу переміщення ребра $(u, v)$ для режиму $m$:

\begin{equation}
  \label{eq:travel_time}
  w_m(u,v) =
  \begin{cases}
    \dfrac{\ell(u,v)}{s_m}
      & \text{якщо режим $m$ дозволений для ребра $(u,v)$,} \\[8pt]
    \dfrac{\ell(u,v)}{s_m \cdot \alpha_{\text{steps}}}
      & \text{якщо } m = \texttt{walk} \text{ та } \texttt{highway} = \texttt{steps},\\[8pt]
    \mathcal{P}
      & \text{інакше (ребро недоступне для режиму $m$),}
  \end{cases}
\end{equation}

\noindent де:
\begin{itemize}
  \item $\ell(u,v) \geq 1$ --- фізична довжина ребра в метрах,
  \item $s_m$ --- прийнята стала швидкість переміщення в метрах за секунду
        (Таблиця~\ref{tab:speeds}),
  \item $\alpha_{\text{steps}} = 0{,}5$ --- коефіцієнт зниження швидкості
        для сходів.
\end{itemize}

\begin{table}[h]
  \centering
  \caption{Прийняті швидкості переміщення за модальним типом.}
  \label{tab:speeds}
  \begin{tabular}{lcc}
    \toprule
    Режим & Швидкість (км/год) & Швидкість (м/с) \\
    \midrule
    Автомобіль (\texttt{drive}) & 30{,}0 & $8{,}\overline{3}$ \\
    Велосипед (\texttt{bike})   & 15{,}0 & $4{,}1\overline{6}$ \\
    Пішки (\texttt{walk})       & 5{,}0  & $1{,}3\overline{8}$ \\
    \bottomrule
  \end{tabular}
\end{table}

Правила модальної доступності, що застосовуються у~\eqref{eq:travel_time},
наведено в Таблиці~\ref{tab:modal_rules}.

\begin{table}[h]
	\centering
	\caption{Правила доступності ребер ($\mathcal{P}$ = непрохідне, інакше --- прохідне).}
	\label{tab:modal_rules}
	\small
	\begin{tabularx}{\textwidth}{lXXX}
		\toprule
		\textbf{Тип ребра} & \textbf{\texttt{drive}} & \textbf{\texttt{bike}} & \textbf{\texttt{walk}} \\
		\midrule
		\texttt{highway} $\in$ \{footway, pedestrian, steps, \dots\}
		& $\mathcal{P}$ & прохідне & прохідне (сходи: $\times 0{,}5$) \\
		\addlinespace % Додає трохи повітря між рядками
		\texttt{motor\_vehicle} $\in$ \{no, private, destination\}
		& $\mathcal{P}$ & прохідне & прохідне \\
		\addlinespace
		\texttt{bicycle} $\in$ \{no, dismount\}
		& $\mathcal{P}$ & $\mathcal{P}$ & прохідне \\
		\addlinespace
		\texttt{highway} = cycleway
		& прохідне & прохідне & прохідне \\
		\addlinespace
		Зворотнє ребро вулиці
		& $\mathcal{P}$ & $\mathcal{P}$ & прохідне \\
		\addlinespace
		\texttt{foot} $\in$ \{no, private\}
		& $\mathcal{P}$ & $\mathcal{P}$ & $\mathcal{P}$ \\
		\bottomrule
	\end{tabularx}
\end{table}

% ─────────────────────────────────────────────────────────────────────────────
\section{Матриця найкоротших відстаней}
% ─────────────────────────────────────────────────────────────────────────────

\subsection{Найкоротші шляхи між усіма парами вузлів (алгоритм Дейкстри)}

Для модального графа $G_m = (V^*, E, w_m)$ та множини локацій
$L = \{0, 1, \ldots, n\}$, де вузол $0$ --- \emph{депо}, а вузли
$1, \ldots, n$ --- точки доставки, \emph{матриця відстаней} визначається як:

\begin{equation}
  \label{eq:matrix}
  d_m(i, j) =
  \begin{cases}
    0 & \text{якщо } i = j, \\
    \displaystyle\min_{\text{шлях } p: i \to j}
      \sum_{(u,v) \in p} w_m(u,v)
      & \text{якщо шлях зі скінченною вартістю існує,} \\
    \mathcal{P} & \text{інакше.}
  \end{cases}
\end{equation}

Мінімізація в~\eqref{eq:matrix} обчислюється алгоритмом Дейкстри для кожної
впорядкованої пари $(i, j)$, де $i \neq j$:

\[
  d_m(i, j) = \operatorname{Dijkstra}(G_m,\, i,\, j,\, w_m).
\]

Побудова повної матриці вимагає $n(n+1)$ запусків алгоритму Дейкстри
($(n+1)$ вузлів, $n(n+1)$ впорядкованих пар).

\subsection{Штрафний сигнал}

Будь-яка пара $(i, j)$ зі значенням $d_m(i, j) \geq \mathcal{P}$
вважається \emph{недосяжною для режиму $m$}. Значення $\mathcal{P} = 10^9$
секунд ($\approx 31{,}7$ року) обрано таким чином, щоб домінувати над будь-яким
реалістичним часом переміщення, залишаючись скінченним числом, --- завдяки
цьому TSP-розв'язувачі можуть будувати коректний (хоча й неоптимальний)
маршрут навіть за наявності недосяжних зупинок.

\subsection{Гібридна маршрутизація «остання миля» (режим автомобіля)}

В режимі автомобіля зупинка $j$ може розташовуватись у пішохідній зоні,
недоступній для автомобіля. У такому разі автомобіль паркується у найближчому
доступному для нього вузлі $v_j^{\text{car}} \in V^*$, а водій проходить
відстань пішки. Скоригована вартість проїзду до зупинки $j$ зі зупинки $i$:

\begin{equation}
  \label{eq:hybrid}
  d_{\text{hybrid}}(i, j) = d_{\text{drive}}(i,\, v_j^{\text{car}})
    + \frac{\ell(v_j^{\text{car}},\, j)}{s_{\text{walk}}},
\end{equation}

\noindent де $\ell(v_j^{\text{car}}, j)$ --- відстань по прямій (формула
Гаверсинуса) від вузла паркування до фактичного місця зупинки.

% ─────────────────────────────────────────────────────────────────────────────
\section{Один транспортний засіб: задача комівояжера}
% ─────────────────────────────────────────────────────────────────────────────

\subsection{Постановка задачі}

Нехай $L = \{0, 1, \ldots, n\}$, де $0$ --- депо. \emph{Маршрут} --- це
перестановка $\pi$ множини $\{1, \ldots, n\}$, що визначає порядок відвідування
зупинок. Загальна вартість кільцевого маршруту:

\begin{equation}
  \label{eq:tsp_cost}
  C(\pi) = d(0,\, \pi_1)
    + \sum_{k=1}^{n-1} d(\pi_k,\, \pi_{k+1})
    + d(\pi_n,\, 0),
\end{equation}

\noindent де $d = d_m$ --- матриця відстаней для обраного режиму $m$.

Задача комівояжера полягає в пошуку перестановки $\pi^*$, що мінімізує
вартість маршруту:

\begin{equation}
  \label{eq:tsp}
  \pi^* = \underset{\pi \in \Pi_n}{\arg\min}\; C(\pi),
\end{equation}

\noindent де $\Pi_n$ --- множина всіх перестановок $n$ зупинок.

Маршрут кодується як \emph{замкнута послідовність}, що починається та
закінчується у депо:
\[
  [0,\; \pi_1,\; \pi_2,\; \ldots,\; \pi_n,\; 0].
\]

\subsection{Методи розв'язання}

Задача комівояжера є NP-важкою в загальному випадку. Система автоматично
обирає евристичний або апроксимаційний алгоритм залежно від кількості
зупинок $n$ (Таблиця~\ref{tab:tsp_methods}).

\begin{table}[h]
  \centering
  \caption{Автоматичний вибір методу розв'язання TSP.}
  \label{tab:tsp_methods}
  \begin{tabular}{cc}
    \toprule
    Кількість зупинок $n$ & Метод \\
    \midrule
    $1$--$2$  & Метод найближчого сусіда (NN) \\
    $3$--$20$ & NN + 2-opt \\
    $21+$     & Генетичний алгоритм (OX) \\
    довільно  & Алгоритм Крістофідеса (явний вибір) \\
    \bottomrule
  \end{tabular}
\end{table}

\subsubsection{Метод найближчого сусіда (NN)}

Починаючи з депо, жадібно додавати найближчу невідвідану зупинку на кожному
кроці:

\begin{equation}
  \label{eq:nn}
  \pi_{k+1} = \underset{j \,\notin\, \{\pi_1, \ldots, \pi_k\}}{\arg\min}\;
    d(\pi_k,\, j), \qquad \pi_0 = 0.
\end{equation}

Евристика NN працює за час $O(n^2)$ і будує коректний маршрут навіть тоді,
коли деякі елементи матриці містять штрафне значення $\mathcal{P}$.

\subsubsection{Локальний пошук 2-opt}

Починаючи з маршруту NN, ітеративно замінювати два несуміжні ребра
$(r_i, r_{i+1})$ та $(r_j, r_{j+1})$, якщо заміна зменшує загальну вартість.
Для маршруту $r = [r_0, r_1, \ldots, r_n, r_0]$ \emph{заміна 2-opt}
на позиціях $i$ та $j$ обертає підпослідовність $r_i, \ldots, r_j$:

\[
  r' = [r_0, \ldots, r_{i-1},\; r_j, r_{j-1}, \ldots, r_i,\; r_{j+1}, \ldots, r_n, r_0].
\]

Заміна приймається, якщо:

\begin{equation}
  \label{eq:2opt}
  C(r') < C(r).
\end{equation}

Процедура повторюється до тих пір, поки не залишиться жодної покращуючої
заміни, або не буде досягнуто максимуму $2\,000$ ітерацій. Складність одного
проходу: $O(n^2)$.

\subsubsection{Апроксимаційний алгоритм Крістофідеса}

Алгоритм Крістофідеса дає $\tfrac{3}{2}$-наближення для метричного TSP.
Він вимагає \emph{повного, зв'язного, неорієнтованого} допоміжного графа
$H = (L, E_H)$, де $E_H$ містить лише ребра з $d(i,j) < \mathcal{P}$.
Якщо $H$ незв'язний (тобто деяка зупинка недосяжна), алгоритм переходить до
методу 2-opt. У застосовному випадку виконуються такі кроки:

\begin{enumerate}[noitemsep]
  \item Побудувати мінімальне кістякове дерево $T$ графа $H$.
  \item Знайти множину $O \subset L$ вузлів непарного ступеня в $T$.
  \item Побудувати мінімальне досконале паросполучення $M$ на підграфі,
        індукованому $O$.
  \item Утворити ейлерів мультиграф $H' = T \cup M$ та виділити ейлерів цикл.
  \item Виключити повторювані вершини, отримавши гамільтонів маршрут.
\end{enumerate}

\subsubsection{Генетичний алгоритм (OX-схрещування)}

Для великих екземплярів ($n \geq 21$) застосовується еволюційний алгоритм.
\emph{Хромосома} --- це перестановка $\pi$ зупинок $\{1,\ldots,n\}$ (без
депо). \emph{Функція пристосованості}:

\begin{equation}
  \label{eq:fitness}
  f(\pi) = \frac{1}{C(\pi) + 1},
\end{equation}

\noindent тобто менша вартість маршруту відповідає більшій пристосованості.

\paragraph{Схрещування за порядком (OX).}
Для батьків $p_1$ та $p_2$ обирається випадковий підрядок $[a, b)$:

\begin{enumerate}[noitemsep]
  \item Скопіювати $p_1[a:b]$ у нащадка на позиції $[a, b)$.
  \item Заповнити решту позицій зліва направо генами з $p_2$ у їх
        початковому порядку, пропускаючи вже розміщені гени.
\end{enumerate}

\paragraph{Мутація обміном.}
З імовірністю $\mu = 0{,}02$ поміняти місцями два випадково обрані гени:
$\pi_i \leftrightarrow \pi_j$.

\paragraph{Генераційний цикл.}
При розмірі популяції $N_{\text{pop}} = 120$ та $G_{\text{max}} = 400$
поколіннях:
\begin{enumerate}[noitemsep]
  \item Відсортувати популяцію за пристосованістю (за спаданням).
  \item Зберегти найкращі $10\%$ як еліту.
  \item Заповнити решту шляхом OX-схрещування випадково обраних особин еліти
        з подальшою мутацією обміном.
\end{enumerate}

Найкраща хромосома після $G_{\text{max}}$ поколінь повертається як
наближений розв'язок.

% ─────────────────────────────────────────────────────────────────────────────
\section{Кілька транспортних засобів: CVRP}
% ─────────────────────────────────────────────────────────────────────────────

\subsection{Постановка задачі}

Введемо позначення:
\begin{itemize}
  \item $S = \{s_1, \ldots, s_n\}$ --- множина точок доставки,
  \item $\mathcal{V} = \{v_1, \ldots, v_K\}$ --- парк з $K$ транспортних засобів,
  \item $m_k \in \mathcal{M}$ --- режим пересування транспортного засобу $v_k$,
  \item $Q_k \in \mathbb{Z}_{>0}$ --- максимальна кількість зупинок для $v_k$,
  \item $d_0 \in V^*$ --- вузол депо.
\end{itemize}

Задача CVRP полягає в пошуку розподілу $\sigma : S \to \mathcal{V}$ та
маршруту $\pi^{(k)}$ для кожного транспортного засобу $v_k$ по призначених
йому зупинках, що мінімізує загальний час переміщення парку:

\begin{equation}
  \label{eq:cvrp}
  \min_{\sigma,\, \{\pi^{(k)}\}} \sum_{k=1}^{K} C_k\!\left(\pi^{(k)}\right)
\end{equation}

\noindent за умов:

\begin{align}
  \sigma(s_i) \neq \varnothing
    &\implies d_{m_{\sigma(s_i)}}(d_0, s_i) < \mathcal{P}
    & &\text{(сумісність режиму)}, \label{eq:cvrp_mode}\\
  \left|\sigma^{-1}(v_k)\right| &\leq Q_k
    & &\forall\, k \in \{1,\ldots,K\}\quad\text{(вантажопідйомність)},\label{eq:cvrp_cap}\\
  \sigma^{-1}(v_k) &\cap \sigma^{-1}(v_j) = \varnothing
    & &\forall\, k \neq j\quad\text{(кожна зупинка --- одному ТЗ)}, \label{eq:cvrp_part}
\end{align}

\noindent де $C_k(\pi^{(k)})$ --- вартість маршруту~\eqref{eq:tsp_cost} для
транспортного засобу $v_k$ з матрицею відстаней $d_{m_k}$.

Зупинки, недосяжні для будь-якого транспортного засобу парку
($d_m(d_0, s_i) \geq \mathcal{P}$ для всіх $m \in \{m_k\}$), виключаються
з оптимізації та повідомляються як попередження.

\subsection{Двофазний алгоритм розв'язання}

Задача CVRP розв'язується шляхом декомпозиції \emph{«спочатку кластеризація,
потім маршрутизація»} у два етапи.

\subsubsection{Фаза 1: Призначення зупинок за модальним типом}

Кожна зупинка $s_i$ призначається тому транспортному режиму $m^*(s_i)$,
транспортні засоби якого мають найбільшу сумарну залишкову місткість серед
сумісних режимів:

\begin{equation}
  \label{eq:mode_assign}
  m^*(s_i) = \underset{m \in \mathcal{M}(s_i)}{\arg\max}\;
    R_m, \qquad
  R_m = \sum_{\substack{k:\, m_k = m}} \bigl(Q_k - |\sigma^{-1}(v_k)|\bigr),
\end{equation}

\noindent де $\mathcal{M}(s_i) = \{m \in \mathcal{M} : d_m(d_0, s_i) < \mathcal{P}\}$
--- множина режимів, сумісних зі зупинкою $s_i$, а $R_m$ --- загальна залишкова
місткість усіх транспортних засобів режиму $m$.

Після призначення режимів зупинки розбиваються на пули за модальним типом:
$\mathcal{S}_m = \{s_i : m^*(s_i) = m\}$.

\subsubsection{Фаза 2: Внутрішньомодальний розподіл за допомогою k-середніх}

У межах кожного модального пулу $\mathcal{S}_m$ зупинки розподіляються між
$K_m = |\{k : m_k = m\}|$ транспортними засобами цього режиму шляхом
географічної кластеризації методом k-середніх.

\paragraph{Кластеризація методом k-середніх.}
Нехай $K_m' = \min(K_m,\, |\mathcal{S}_m|)$. Застосувати метод k-середніх до
географічних координат зупинок пулу $\mathcal{S}_m$:

\begin{equation}
  \label{eq:kmeans}
  \min_{\{C_1,\ldots,C_{K_m'}\}}
    \sum_{j=1}^{K_m'} \sum_{s_i \in C_j}
    \left\| \begin{pmatrix} \varphi(s_i) \\ \lambda(s_i) \end{pmatrix}
           - \boldsymbol{\mu}_j \right\|^2,
\end{equation}

\noindent де $\varphi(s_i)$ та $\lambda(s_i)$ --- широта та довгота зупинки
$s_i$, $\boldsymbol{\mu}_j$ --- центроїд кластера $j$, а $\|\cdot\|$ ---
евклідова норма.

\paragraph{Призначення кластерів транспортним засобам за азимутом.}
Для географічно узгодженого розподілу кластерів центроїди впорядковуються
за \emph{компасним азимутом} відносно депо. Азимут від точки
$(\varphi_1, \lambda_1)$ до точки $(\varphi_2, \lambda_2)$:

\begin{equation}
  \label{eq:bearing}
  \theta = \left(
    \arctan2\!\left(
      \sin(\Delta\lambda)\cos\varphi_2,\;
      \cos\varphi_1 \sin\varphi_2 - \sin\varphi_1 \cos\varphi_2 \cos(\Delta\lambda)
    \right) \cdot \frac{180}{\pi} + 360
  \right) \bmod 360,
\end{equation}

\noindent де $\Delta\lambda = \lambda_2 - \lambda_1$. Кластери призначаються
транспортним засобам у порядку зростання азимута.

\paragraph{Перебалансування місткості.}
Метод k-середніх не враховує місткість транспортних засобів. Після початкового
призначення будь-який транспортний засіб $v_k$, що перевищив свою місткість
$Q_k$, передає надлишкові зупинки наступному транспортному засобу з вільною
місткістю:

\begin{equation}
  \label{eq:rebalance}
  \text{поки } |\sigma^{-1}(v_k)| > Q_k :
  \quad\text{видалити останню зупинку з } \sigma^{-1}(v_k),
  \quad\text{додати до першого } v_j \text{, де } |\sigma^{-1}(v_j)| < Q_j.
\end{equation}

Кандидати перебираються спочатку у прямому напрямку ($j > k$), потім у
зворотному ($j < k$), щоб уникнути хибної помилки переповнення для останнього
транспортного засобу за азимутом. Якщо жодного транспортного засобу з
вільною місткістю не знайдено, генерується виключення \texttt{ValueError}.

\subsubsection{TSP для кожного транспортного засобу}

Після визначення розподілу $\sigma$ кожен транспортний засіб $v_k$ зі
$|\sigma^{-1}(v_k)| \geq 1$ незалежно розв'язує задачу
TSP~\eqref{eq:tsp}--\eqref{eq:tsp_cost} на своєму кластері, використовуючи
модальну матрицю відстаней $d_{m_k}$ та автоматично обраний метод з
Таблиці~\ref{tab:tsp_methods}.

Загальна вартість парку декомпозується як:

\begin{equation}
  \label{eq:fleet_cost}
  C_{\text{fleet}} = \sum_{k=1}^{K} C_k(\pi^{(k)}).
\end{equation}

% ─────────────────────────────────────────────────────────────────────────────
\section{Відновлення маршруту}
% ─────────────────────────────────────────────────────────────────────────────

Розв'язок TSP $\pi^*$ є послідовністю вузлів вищого рівня
$[0, \pi_1^*, \ldots, \pi_n^*, 0]$. Для відображення маршруту на карті кожна
послідовна пара $(r_k, r_{k+1})$ розгортається у повну послідовність
проміжних OSM-вузлів дороги за допомогою алгоритму Дейкстри:

\begin{equation}
  \label{eq:reconstruct}
  \text{шлях}(r_k, r_{k+1}) = \underset{p:\, r_k \to r_{k+1}}{\arg\min}
    \sum_{(u,v) \in p} w_m(u, v).
\end{equation}

Конкатенація всіх розгорнутих підшляхів (зі видаленням дублікатів спільних
граничних вузлів) утворює \emph{повний ламаний маршрут}, що відображається
на карті Folium.

% ─────────────────────────────────────────────────────────────────────────────
\section{Зведення позначень}
% ─────────────────────────────────────────────────────────────────────────────

\begin{table}[h]
  \centering
  \caption{Зведення математичних позначень.}
  \begin{tabular}{cl}
    \toprule
    Символ & Опис \\
    \midrule
    $G = (V, E, w)$ & Орієнтований зважений мультиграф дорожньої мережі \\
    $G^*$ & Найбільша сильно зв'язна компонента графа $G$ \\
    $G_m$ & Модальна копія $G^*$ для режиму $m$ \\
    $w_m(u, v)$ & Ваговий коефіцієнт часу проходження ребра $(u, v)$ у режимі $m$ (с) \\
    $\mathcal{P} = 10^9$ & Штрафний сигнал непрохідності / недосяжності (с) \\
    $\ell(u, v)$ & Фізична довжина ребра $(u, v)$ в метрах \\
    $s_m$ & Швидкість переміщення для режиму $m$ в метрах за секунду \\
    $d_m(i, j)$ & Час найкоротшого шляху від $i$ до $j$ у режимі $m$ \\
    $n$ & Кількість точок доставки (без депо) \\
    $L = \{0, 1, \ldots, n\}$ & Множина локацій; вузол $0$ --- депо \\
    $\pi$ & Перестановка зупинок (маршрут TSP, без депо) \\
    $C(\pi)$ & Загальна вартість кільцевого маршруту \\
    $\pi^*$ & Оптимальний (або наближено оптимальний) маршрут \\
    $K$ & Кількість транспортних засобів у парку \\
    $\mathcal{V} = \{v_1, \ldots, v_K\}$ & Парк транспортних засобів \\
    $m_k$ & Режим пересування транспортного засобу $v_k$ \\
    $Q_k$ & Максимальна кількість зупинок для транспортного засобу $v_k$ \\
    $\sigma : S \to \mathcal{V}$ & Функція призначення зупинок транспортним засобам \\
    $\mathcal{S}_m$ & Множина зупинок, призначених до модального пулу $m$ \\
    $C_j$ & $j$-й кластер зупинок методу k-середніх \\
    $\boldsymbol{\mu}_j$ & Центроїд кластера $C_j$ \\
    $\theta$ & Компасний азимут (градуси, $[0, 360)$) \\
    $f(\pi)$ & Пристосованість генетичного алгоритму (обернена вартість) \\
    \bottomrule
  \end{tabular}
\end{table}

\end{document}
